\documentclass[UTF8,fontset=fandol,a4paper,12pt]{ctexart}
\usepackage[margin=2.5cm]{geometry}
\usepackage{amsmath,booktabs,graphicx,listings}
\usepackage[backend=biber,style=numeric]{biblatex}
\addbibresource{refs.bib}
\usepackage[hidelinks]{hyperref}
\title{系统实验 01：观察数组的存储布局}
\author{姓名：\underline{\hspace{3cm}}}
\date{\today}
\begin{document}
\maketitle
\begin{abstract}
用 C 程序输出数组元素的地址和元素大小，用相邻地址的差检验数组布局。
表中的数值是演算示例，提交时请替换为自己的测量。
\end{abstract}
\section{问题与预测}
记元素大小为 $s$ 字节，首元素的字节地址为 $A_0$。
连续数组中，下标为 $i$ 的元素地址应满足式~\eqref{eq:address}。
\begin{equation}\label{eq:address}
  A_i=A_0+i\,s .
\end{equation}
有关机器级数据表示的背景参见文献~\cite{csapp3}。
\section{方法}
记录操作系统、编译器版本与编译参数，再运行以下程序。
\begin{lstlisting}[language=C,basicstyle=\ttfamily\small,frame=single]
#include <stdio.h>
int main(void) {
    int a[3] = {10, 20, 30};
    printf("size = %zu bytes\n", sizeof a[0]);
    for (size_t i = 0; i < 3; ++i)
        printf("%zu: %p\n", i, (void *)&a[i]);
    return 0;
}
\end{lstlisting}
\section{结果}
图~\ref{fig:array} 表示元素顺序，表~\ref{tab:offset} 给出示例偏移。
\begin{figure}[htbp]
  \centering
  \fbox{\strut $a[0]$}\fbox{\strut $a[1]$}\fbox{\strut $a[2]$}
  \caption{数组元素的排列示意}\label{fig:array}
\end{figure}
\begin{table}[htbp]
  \centering
  \caption{当元素大小为 4 字节时的示例偏移}\label{tab:offset}
  \begin{tabular}{rrr}
    \toprule
    下标 & 值 & 相对首地址的偏移（字节）\\
    \midrule
    0 & 10 & 0\\
    1 & 20 & 4\\
    2 & 30 & 8\\
    \bottomrule
  \end{tabular}
\end{table}
\section{讨论}
请填入实际测得的元素大小、三条地址和两次相邻地址差。
比较测量与预测，说明换成 \texttt{double} 后应如何修改预测。
\printbibliography[title={参考文献}]
\end{document}
